\section{Method}
\label{sec:method}

\subsection{Overview}
\label{sec:method:overview}
YouDrive keeps a VLA backbone as a scene encoder and attaches a style-conditioned
trajectory decoder. Multi-view images, ego state, and the navigation command are
encoded into a sequence of content tokens. A persona is added to these tokens as a
style direction with a single coefficient, and a flow-matching decoder samples
several feasible trajectories from the resulting content and returns one.
\Cref{fig:pipeline} shows the pipeline. The design follows the two properties the
average-driver problem calls for: style is an explicit and composable axis, and the
decoder can commit to a mode rather than average over several.

\begin{figure}[t]
    \centering
    \fbox{\parbox[c][3.0cm][c]{0.9\linewidth}{\centering\small
    \TODO{Figure 2: pipeline. VLA backbone produces content tokens; a persona task
    vector is interpolated in with coefficient $\alpha$; the flow-matching decoder
    samples $N{=}16$ candidates; one is returned.}}}
    \caption{YouDrive pipeline. The VLA backbone provides content tokens, style is
    a task-vector interpolation applied to those tokens, and a flow-matching decoder
    samples several feasible candidates before returning one.}
    \label{fig:pipeline}
\end{figure}

\subsection{Style as a task vector}
\label{sec:method:style}
We represent a driving style as a direction that can be added to the base policy.
For each teacher planner we distill a low-rank adapter on top of the frozen
backbone; its weight increment $\Delta W$ is the persona's task vector in the sense
of task arithmetic, a reusable edit that moves the base policy toward the teacher's
behavior. Since the backbone is fixed and only the adapter is trained, all personas
are defined against a common reference and can be scaled and combined.

We expose the task vector through one user coefficient rather than by editing
weights at run time. Let $h_{\mathrm{base}}$ be the content produced by the
backbone and $h_{\mathrm{styled}}$ the content after applying a persona adapter. We
condition the decoder on
\begin{equation}
    c = h_{\mathrm{base}} + \alpha\,(h_{\mathrm{styled}} - h_{\mathrm{base}}),
    \qquad \alpha \in [0,1],
    \label{eq:interp}
\end{equation}
and combine several personas by
$c = h_{\mathrm{base}} + \sum_i w_i\,(h^{(i)}_{\mathrm{styled}} - h_{\mathrm{base}})$.
This is a first-order realization of task-vector scaling and addition in the
content space of the VLA: $\alpha$ moves from the average driver at $\alpha{=}0$ to
the persona at $\alpha{=}1$, and the weights $w_i$ blend personas. In
\cref{sec:exp:vector} we check that the induced style shifts are close to linear in
$\alpha$, additive under blending, and reversed under negation, which is the
behavior expected of a task vector.

Working in the VLA content space is part of the method, not an implementation
detail. Task-vector control has been reported in language and vision-language
representations, where directions align with interpretable behavior. A geometric
bird's-eye or image encoder organizes its features around scene geometry rather
than around how to act, so a style direction defined there need not be linear or
composable. VLA content jointly encodes what the scene contains and what the agent
is asked to do, including the navigation command and interaction cues, which is the
setting in which a behavioral direction is well defined. \Cref{sec:exp:ablate}
compares the same personas defined in a matched geometric feature space and finds
weaker control.

\subsection{Flow-matching trajectory decoder}
\label{sec:method:fm}
The backbone we build on decodes a trajectory by selecting one token sequence from
a discrete action codebook, which commits to a single mode before the trajectory is
formed. We replace this decoder with a conditional flow-matching head that regresses
a velocity field from a noise sample to a trajectory~\cite{rectflow,flowmatching}
using a DiT/AdaLN design~\cite{dit}, conditioned on the content $c$. The head
decodes in a control space of longitudinal acceleration and yaw rate under a
kinematic unicycle model, so every integrated sample is dynamically feasible. At
inference we integrate the flow with $K$ steps and draw $N{=}16$ samples in
parallel, producing a set of candidate trajectories instead of a point estimate.
This change is what lets a persona express itself in multi-modal scenes: distinct
maneuvers appear as distinct samples rather than being averaged into one.

\subsection{From samples to a trajectory}
\label{sec:method:select}
The decoder produces $N$ candidates per frame and the planner returns one. We score
the candidates with a lightweight scorer and keep the highest-scoring one; a learned
scorer fits the same interface. The style of the returned trajectory is a property
of the sampled set, so style control does not depend on the choice of scorer.
